\documentclass[12pt,a4paper]{article}

\usepackage[T1]{fontenc}
\usepackage[utf8]{inputenc}
\usepackage{geometry}
\usepackage{amsmath,amssymb,amsthm}
\usepackage{mathtools}
\usepackage{array}
\usepackage{booktabs}
\usepackage{colortbl}
\usepackage{enumitem}
\usepackage{xcolor}
\usepackage{fancyhdr}
\usepackage{titlesec}
\usepackage{tcolorbox}
\usepackage{listings}
\usepackage{parskip}
\usepackage{longtable}
\usepackage{tikz}

\tcbuselibrary{skins,breakable}

\geometry{
  top=2.5cm, bottom=2.5cm,
  left=3.0cm, right=2.5cm,
  headheight=15pt
}

\definecolor{bleuTitre}{RGB}{31,56,100}
\definecolor{bleuSection}{RGB}{31,92,153}
\definecolor{bleuSub}{RGB}{46,117,182}
\definecolor{grisNote}{RGB}{245,244,240}
\definecolor{codebg}{RGB}{248,246,240}
\definecolor{codecomment}{RGB}{100,100,90}
\definecolor{vertOk}{RGB}{30,120,50}
\definecolor{rougeAlert}{RGB}{180,40,20}
\definecolor{grisLigne}{RGB}{235,235,230}

\pagestyle{fancy}
\fancyhf{}
\renewcommand{\headrulewidth}{0.4pt}
\fancyhead[L]{\small\textit{Validation et reconstruction 3D}}
\fancyfoot[C]{\thepage}

\titleformat{\section}
  {\large\bfseries\color{bleuTitre}}
  {\thesection.}{0.6em}{}[\vspace{-4pt}\rule{\linewidth}{0.6pt}]
\titleformat{\subsection}
  {\normalsize\bfseries\color{bleuSection}}
  {\thesubsection}{0.6em}{}
\titleformat{\subsubsection}
  {\normalsize\itshape\bfseries}
  {\thesubsubsection}{0.6em}{}

\theoremstyle{definition}
\newtheorem{definition}{D\'efinition}[section]
\newtheorem{propriete}{Propri\'et\'e}[section]
\theoremstyle{remark}
\newtheorem{remarque}{Remarque}[section]

\newtcolorbox{notebox}[1][]{
  enhanced, breakable,
  colback=grisNote,
  colframe=bleuSection,
  leftrule=4pt, toprule=0.4pt, bottomrule=0.4pt, rightrule=0.4pt,
  arc=2pt,
  fonttitle=\bfseries\small\color{bleuSection},
  title={#1},
  before upper={\setlength{\parskip}{4pt}},
}

\newtcolorbox{resultatbox}[1][]{
  enhanced, breakable,
  colback=vertOk!5,
  colframe=vertOk!70!black,
  leftrule=4pt, toprule=0.4pt, bottomrule=0.4pt, rightrule=0.4pt,
  arc=2pt,
  fonttitle=\bfseries\small\color{vertOk!70!black},
  title={#1},
  before upper={\setlength{\parskip}{4pt}},
}

\newtcolorbox{alertbox}[1][]{
  enhanced, breakable,
  colback=rougeAlert!5,
  colframe=rougeAlert!70!black,
  leftrule=4pt, toprule=0.4pt, bottomrule=0.4pt, rightrule=0.4pt,
  arc=2pt,
  fonttitle=\bfseries\small\color{rougeAlert!70!black},
  title={#1},
  before upper={\setlength{\parskip}{4pt}},
}

\lstdefinestyle{benzai}{
  backgroundcolor=\color{codebg},
  basicstyle=\ttfamily\small,
  breaklines=true,
  frame=single,
  framesep=6pt,
  rulecolor=\color{bleuSection!40},
  xleftmargin=12pt,
  xrightmargin=6pt,
  commentstyle=\color{codecomment}\itshape,
  morecomment=[l]{\#},
  showstringspaces=false,
}

\lstdefinestyle{python}{
  backgroundcolor=\color{codebg},
  basicstyle=\ttfamily\small,
  breaklines=true,
  frame=single,
  framesep=6pt,
  rulecolor=\color{bleuSection!40},
  xleftmargin=12pt,
  xrightmargin=6pt,
  commentstyle=\color{codecomment}\itshape,
  keywordstyle=\color{bleuSection}\bfseries,
  stringstyle=\color{vertOk!80!black},
  language=Python,
  showstringspaces=false,
  morekeywords={self,True,False,None},
}

\newcommand{\B}{\mathcal{B}}
\newcommand{\GD}{G_{\!D}}
\newcommand{\VD}{V_{\!D}}
\newcommand{\ED}{E_{\!D}}
\newcommand{\xv}{x_v}
\newcommand{\pat}[1]{\mathtt{#1}}

% ────────────────────────────────────────────────────────────────────
\begin{document}

\begin{center}
  {\LARGE\bfseries\color{bleuTitre}
    Validation des solutions CSP~:\\[4pt]
    Reconstruction 3D des mol\'ecules}\\[6pt]
  {\large\itshape
    Pipeline de reconstruction g\'eom\'etrique\\
    et r\'esultats exp\'erimentaux}\\[18pt]
  {\small Document de travail --- Stage M2 IMD --- Avril 2026}\\[4pt]
  {\small\color{gray}\today}
\end{center}

\vspace{8pt}
\rule{\linewidth}{1pt}
\vspace{4pt}

\tableofcontents

\vspace{10pt}
\rule{\linewidth}{0.4pt}


% ════════════════════════════════════════════════════════════════════
\section{Probl\'ematique}
% ════════════════════════════════════════════════════════════════════

\subsection{Contexte}

Le solveur CSP d\'ecrit dans le document pr\'ec\'edent
(\textit{Mod\'elisation du probl\`eme comme un CSP}, v6.0) \'enum\`ere
les assignations de tailles $\{5, 6, 7\}$ aux hexagones d'un
benz\'eno\"ide~$\B$, sous les certaines contraintes.

Chaque solution est un dictionnaire $\{v \mapsto \text{taille}_v\}$
o\`u $v$ parcourt les sommets du graphe dual~$\GD$. Mais une solution
CSP n'est qu'une \textbf{assignation abstraite}~: elle ne dit rien sur la
g\'eom\'etrie 3D de la mol\'ecule r\'esultante. Pour valider chimiquement
une solution, il faut~:
\begin{enumerate}[nosep]
  \item Reconstruire la mol\'ecule compl\`ete (atomes de carbone, liaisons C--C, hydrog\`enes).
  \item Lui attribuer des coordonn\'ees 3D raisonnables.
  \item Optimiser la g\'eom\'etrie avec xTB (GFN2-xTB).
  \item Tester la planarit\'e globale (ACP, seuil de 10$^\circ$).
\end{enumerate}

\subsection{Probl\`eme de la reconstruction}

Le benz\'eno\"ide original poss\`ede une g\'eom\'etrie parfaitement
d\'efinie~: chaque sommet a des coordonn\'ees hexagonales $(h_x, h_y)$
convertibles en cart\'esiennes~:
\[
  X = d_{\text{CC}} \cdot \bigl(h_x + h_y / 2\bigr), \qquad
  Y = d_{\text{CC}} \cdot h_y \cdot \frac{\sqrt{3}}{2}, \qquad
  Z = 0
\]
avec $d_{\text{CC}} = 1{,}42$~\AA.

Cependant, quand un hexagone (angles internes de $120^\circ$)
est remplac\'e par un pentagone ($108^\circ$) ou un heptagone
($\approx 128{,}57^\circ$), les coordonn\'ees hexagonales ne sont
plus valides. On ne peut pas simplement retirer ou ajouter un carbone
sur la grille hexagonale et obtenir une g\'eom\'etrie correcte.

\begin{notebox}[Premi\`ere tentative (\'echec)]
Une premi\`ere approche na\"ive consistait \`a~:
\begin{itemize}[nosep]
  \item Placer tous les carbones sur la grille hexagonale.
  \item Pour un pentagone~: retirer un carbone du bloc libre et reconnecter les voisins.
  \item Pour un heptagone~: couper une ar\^ete libre et ins\'erer un nouveau carbone au milieu, d\'ecal\'e perpendiculairement.
\end{itemize}
Cette approche produit des topologies correctes mais une g\'eom\'etrie
bris\'ee~: les longueurs de liaison et les angles ne correspondent pas
aux polygones r\'eguliers, et xTB ne parvient pas \`a corriger des
distorsions aussi importantes.
\end{notebox}


% ════════════════════════════════════════════════════════════════════
\section{Approche adopt\'ee~: topologie d'abord, g\'eom\'etrie ensuite}
% ════════════════════════════════════════════════════════════════════

L'id\'ee cl\'e est de \textbf{s\'eparer} la reconstruction en deux phases
ind\'ependantes~:
\begin{description}[leftmargin=2.5cm,style=nextline]
  \item[Phase~A --- Topologie]
    Construire le graphe mol\'eculaire correct (quels atomes existent,
    quelles liaisons, quels cycles) en modifiant le benz\'eno\"ide
    d'origine.
  \item[Phase~B --- G\'eom\'etrie]
    Attribuer des coordonn\'ees 2D \`a chaque atome en pla\c{c}ant les
    cycles un par un comme des \textbf{polygones r\'eguliers}, via un
    parcours BFS du graphe dual.
\end{description}

La topologie doit \^etre \textbf{exacte}~: une erreur d'un seul atome
ou d'une seule liaison produit une mol\'ecule chimiquement diff\'erente.
La g\'eom\'etrie, en revanche, peut \^etre \textbf{approximative}~: xTB
corrige les positions atomiques lors de l'optimisation, tant que la
structure initiale est raisonnablement proche de l'\'equilibre.

\subsection{Invariant fondamental}

\begin{propriete}[Conservation des c\^ot\'es partag\'es]
\label{prop:invariant}
Soit un hexagone~$v$ de pattern $\sigma(v) = (\sigma_0, \ldots, \sigma_5)$,
o\`u $\sigma_i = 1$ si le c\^ot\'e~$i$ est partag\'e avec un voisin et
$\sigma_i = 0$ sinon.

Les modifications topologiques (passage \`a un pentagone ou un heptagone)
n'op\`erent que sur le \textbf{bloc libre} (c\^ot\'es o\`u $\sigma_i = 0$).
Par cons\'equent, les c\^ot\'es partag\'es et leurs sommets ne sont
\textbf{jamais affect\'es}.
\end{propriete}

Cet invariant simplifie consid\'erablement l'algorithme~:
\begin{itemize}[nosep]
  \item Les sommets de l'ar\^ete partag\'ee entre deux hexagones sont
        toujours les m\^emes, avant et apr\`es modification.
  \item On peut utiliser ces sommets comme \textbf{ancrage g\'eom\'etrique}
        pour placer les cycles adjacents.
  \item L'indexation $\texttt{edge\_positions}[(u,v)]$ du parseur Benzai
        reste valide apr\`es modification.
\end{itemize}


% ════════════════════════════════════════════════════════════════════
\section{D\'etail de l'algorithme}
% ════════════════════════════════════════════════════════════════════

\subsection{Phase A~: reconstruction topologique}

\subsubsection{Hexagone inchang\'e ($x_v = 6$)}

Le cycle conserve ses 6~sommets et 6~ar\^etes. Aucune modification.

\subsubsection{Pentagone ($x_v = 5$)~--- retrait d'un sommet}

\begin{enumerate}[nosep]
  \item Identifier le \textbf{bloc libre}~: les c\^ot\'es cons\'ecutifs
        o\`u $\sigma_i = 0$.
  \item Identifier les \textbf{sommets int\'erieurs} au bloc libre~: un
        sommet~$h_i$ est int\'erieur si ses deux c\^ot\'es incidents dans
        le cycle (c\^ot\'es $i{-}1$ et~$i$) sont tous deux libres~:
        $\sigma_{(i-1) \bmod 6} = 0$ et $\sigma_i = 0$.
  \item Choisir le sommet int\'erieur du \textbf{milieu} du bloc.
  \item Le retirer du graphe~: supprimer ses deux liaisons incidentes dans
        le cycle et ajouter une liaison directe entre ses deux ex-voisins.
\end{enumerate}

\begin{notebox}[Justification de la s\'ecurit\'e du retrait]
Un sommet int\'erieur au bloc libre n'appartient \`a aucun autre hexagone
du benz\'eno\"ide. En effet, dans le r\'eseau hexagonal, un sommet partag\'e
entre deux hexagones se trouve n\'ecessairement sur un c\^ot\'e partag\'e
($\sigma_i = 1$). Si les deux c\^ot\'es incidents d'un sommet sont libres,
les hexagones voisins correspondants n'existent pas dans le benz\'eno\"ide.
Le retrait est donc \textbf{topologiquement s\^ur}.
\end{notebox}

\subsubsection{Heptagone ($x_v = 7$)~--- insertion d'un sommet}

\begin{enumerate}[nosep]
  \item Identifier le bloc libre.
  \item Choisir le \textbf{c\^ot\'e libre} du milieu du bloc.
  \item Couper l'ar\^ete correspondante et ins\'erer un nouveau sommet
        entre ses deux extr\'emit\'es.
\end{enumerate}

Un c\^ot\'e libre n'est partag\'e avec aucun voisin~: l'insertion ne
perturbe la topologie d'aucun cycle adjacent.

\subsubsection{Construction des ensembles globaux}

Apr\`es les modifications, on dispose de~:
\begin{itemize}[nosep]
  \item \texttt{vertex\_set}~: l'ensemble des labels de carbones
        (originaux survivants $+$ nouveaux sommets ins\'er\'es).
  \item \texttt{bond\_set}~: l'ensemble des liaisons C--C
        (originales $-$ supprim\'ees $+$ ajout\'ees).
  \item \texttt{cycle\_vertices[$v$]}~: pour chaque hexagone~$v$, la liste
        ordonn\'ee des sommets du cycle modifi\'e (de taille $x_v$).
\end{itemize}


\subsection{Phase B~: placement g\'eom\'etrique par BFS}

\subsubsection{Principe}

Chaque cycle est plac\'e comme un \textbf{polygone r\'egulier} de
$n = x_v$~c\^ot\'es, avec une longueur de c\^ot\'e $d_{\text{CC}} = 1{,}42$~\AA.
Les formules de base sont~:
\[
  R(n) = \frac{d_{\text{CC}}}{2\sin(\pi/n)}
  \qquad\text{(rayon circonscrit)},\qquad
  a(n) = \frac{d_{\text{CC}}}{2\tan(\pi/n)}
  \qquad\text{(apoth\`eme)}.
\]

\subsubsection{Parcours BFS du graphe dual}

\begin{enumerate}[nosep]
  \item \textbf{Choix de la racine}~: le sommet du dual de degr\'e
        maximal, avec pr\'ef\'erence pour un hexagone inchang\'e
        ($x_v = 6$), qui offre un ancrage g\'eom\'etrique stable.
  \item \textbf{Placement de la racine}~: polygone r\'egulier centr\'e
        en $(0, 0)$, premier sommet orient\'e vers le haut.
  \item \textbf{BFS}~: pour chaque cycle~$v$ d\'ecouvert avec parent~$p$~:
        \begin{enumerate}[nosep,label=(\alph*)]
          \item R\'ecup\'erer l'ar\^ete partag\'ee via
                $\texttt{edge\_positions}[(p, v)]$. Les deux sommets de
                cette ar\^ete sont d\'ej\`a plac\'es.
          \item Calculer le centre du nouveau cycle~:
                milieu de l'ar\^ete $+$ $a(x_v)$ dans la direction
                perpendiculaire, orient\'ee \`a l'oppos\'e du centre de~$p$.
          \item Placer les sommets comme polygone r\'egulier autour de ce
                centre, ancr\'e sur le premier sommet partag\'e.
          \item \textbf{R\`egle de snap}~: si un sommet est d\'ej\`a
                plac\'e (partag\'e avec un autre cycle d\'ej\`a trait\'e),
                conserver ses coordonn\'ees existantes.
        \end{enumerate}
\end{enumerate}

\begin{notebox}[Gestion de l'adjacence multiple]
Un cycle peut partager des ar\^etes avec plusieurs cycles d\'ej\`a plac\'es
(pas seulement son parent BFS). La r\`egle de snap (\'etape~3d) r\'esout
ce cas~: les sommets partag\'es gardent leurs coordonn\'ees, et le cycle
n'est donc pas un polygone parfaitement r\'egulier. Les distorsions
r\'esultantes sont mineures et corrig\'ees par l'optimisation xTB.
\end{notebox}

\subsubsection{D\'etermination du sens de rotation}

L'ar\^ete partag\'ee est travers\'ee en sens inverse par le parent et
l'enfant (propri\'et\'e des cycles adjacents dans un graphe planaire).
Le sens de rotation pour le placement est donc d\'etermin\'e depuis la
\textbf{perspective de l'enfant}~: on calcule la direction angulaire de
$\texttt{child\_v1}$ vers $\texttt{child\_v2}$ vus depuis le centre de
l'enfant, et on place les sommets restants en poursuivant dans cette
direction.


\subsection{Phase C~: assemblage et export}

Une fois la topologie et les coordonn\'ees \'etablies~:
\begin{enumerate}[nosep]
  \item Construire un \texttt{MolecularGraph} avec tous les carbones
        (coordonn\'ees 2D, $z = 0$) et toutes les liaisons C--C.
  \item Enregistrer les cycles (n\'ecessaires pour l'\'etape suivante).
  \item Appliquer le \texttt{ValenceSolver}~: placement des doubles
        liaisons par couplage maximum pond\'er\'e (algorithme de Blossom),
        puis placement des hydrog\`enes selon la proc\'edure du chimiste.
  \item Exporter en format XYZ pour xTB.
\end{enumerate}


% ════════════════════════════════════════════════════════════════════
\section{Architecture logicielle}
% ════════════════════════════════════════════════════════════════════

Le pipeline de reconstruction est impl\'ement\'e dans un package d\'edi\'e
\texttt{csp\_solver/reconstruction/}, s\'epar\'e du solveur CSP
lui-m\^eme.

\begin{center}
\begin{tabular}{lll}
\toprule
\textbf{Module} & \textbf{R\^ole} & \textbf{Phase} \\
\midrule
\texttt{topology.py}  & Modifications topologiques (pentagone/heptagone) & A \\
\texttt{placement.py}  & Placement g\'eom\'etrique BFS (polygones r\'eguliers) & B \\
\texttt{assembler.py}  & Construction \texttt{MolecularGraph} + \texttt{ValenceSolver} + export & C \\
\texttt{pipeline.py}   & Orchestration + validation xTB + planarit\'e & --- \\
\texttt{\_\_init\_\_.py}  & R\'eexport de l'API publique & --- \\
\bottomrule
\end{tabular}
\end{center}

Chaque module a une \textbf{responsabilit\'e unique}, ce qui facilite
le d\'ebogage, la maintenance et les \'evolutions futures. Les
d\'ependances vers le g\'en\'erateur existant
(\texttt{non\_benzenoid\_generator/}) sont limit\'ees aux structures de
donn\'ees (\texttt{MolecularGraph}, \texttt{Cycle}) et au
\texttt{ValenceSolver}.

\begin{notebox}[Invocation]
\begin{lstlisting}[style=benzai]
cd csp_solver/
python main.py data/first.graph --validate
python main.py data/second.graph --validate
\end{lstlisting}
Le flag \texttt{--validate} active le pipeline complet~: reconstruction
$\to$ export XYZ $\to$ optimisation xTB $\to$ test de planarit\'e.
\end{notebox}


% ════════════════════════════════════════════════════════════════════
\section{R\'esultats exp\'erimentaux}
% ════════════════════════════════════════════════════════════════════

\subsection{Test 1~: \texttt{first.graph} --- 3 hexagones lin\'eaires}

Le benz\'eno\"ide \texttt{first.graph} est un alignement lin\'eaire de
3~hexagones ($h = 3$, $|V| = 14$, $|\ED| = 2$). Le solveur CSP produit
\textbf{2~solutions} (apr\`es rupture de sym\'etrie par lex-leader)~:

\begin{center}
\begin{tabular}{clcc}
\toprule
\textbf{\#} & \textbf{Assignation} & \textbf{Carbones} & \textbf{Planarit\'e} \\
\midrule
1 & $v_0 = 5,\; v_1 = 6,\; v_2 = 7$ & 14 & \textcolor{vertOk}{\textbf{PLAN}} \\
2 & $v_0 = 6,\; v_1 = 6,\; v_2 = 6$ & 14 & \textcolor{vertOk}{\textbf{PLAN}} \\
\bottomrule
\end{tabular}
\end{center}

\begin{resultatbox}[R\'esultat --- \texttt{first.graph}]
Les 2~solutions sont reconstruites, optimis\'ees par xTB, et
valid\'ees comme planes. La solution~1 (pentagone--hexagone--heptagone)
correspond \`a une structure de type azul\'enique fusionn\'ee.
La solution~2 est le benz\'eno\"ide original.
\end{resultatbox}


\subsection{Test 2~: \texttt{second.graph} --- 16 hexagones}

Le benz\'eno\"ide \texttt{second.graph} est une structure plus complexe
($h = 16$, $|V| = 60$, $|\ED| = 19$). Le solveur CSP produit
\textbf{41~solutions} distinctes.

\begin{center}
\begin{tabular}{lc}
\toprule
\textbf{M\'etrique} & \textbf{Valeur} \\
\midrule
Solutions test\'ees & 41 \\
Planes (valides, angle $\leq 10^\circ$) & \textcolor{rougeAlert}{\textbf{0}} \\
Non planes (angle $> 10^\circ$) & \textcolor{rougeAlert}{\textbf{41}} \\
\bottomrule
\end{tabular}
\end{center}

\begin{alertbox}[R\'esultat --- \texttt{second.graph}]
Aucune des 41~solutions n'est plane apr\`es optimisation xTB. Les
mol\'ecules reconstruites ont une g\'eom\'etrie 3D qui d\'evie
significativement du plan, avec des angles sup\'erieurs au seuil de
$10^\circ$.
\end{alertbox}


% ════════════════════════════════════════════════════════════════════
\section{Discussion et perspectives}
% ════════════════════════════════════════════════════════════════════

\subsection{Fiabilit\'e de la reconstruction}

Le pipeline de reconstruction a \'et\'e valid\'e sur le cas
\texttt{first.graph}~:
\begin{itemize}[nosep]
  \item La solution tout-hexagonale ($v_0 = v_1 = v_2 = 6$) reproduit
        le benz\'eno\"ide original et passe le test de planarit\'e.
  \item La solution mixte ($5$--$6$--$7$) produit une mol\'ecule
        topologiquement correcte, g\'eom\'etriquement raisonnable, et
        valid\'ee comme plane par xTB.
\end{itemize}

L'inspection visuelle des fichiers XYZ g\'en\'er\'es pour
\texttt{second.graph} confirme que les mol\'ecules sont
topologiquement correctes~: les cycles ont les bonnes tailles, les
liaisons sont coh\'erentes, et les hydrog\`enes sont bien plac\'es.
Le pipeline de reconstruction fonctionne donc correctement.

\subsection{Suffisance des contraintes CSP}

Le fait que les 41~solutions de \texttt{second.graph} soient
\textbf{toutes non planes} soul\`eve une question fondamentale~:

\begin{alertbox}[Question ouverte]
Les contraintes CSP actuelles (conservation du carbone~C1, blocs de
z\'eros~C2, voisinage admissible~C3) sont-elles \textbf{suffisantes}
pour garantir l'existence d'un plongement plan de la mol\'ecule
r\'esultante~?
\end{alertbox}

Nos contraintes raisonnent uniquement sur le \textbf{voisinage local}~:
la table $\mathcal{T}(n)$ v\'erifie qu'un cycle central et ses voisins
imm\'ediats forment une configuration g\'eom\'etriquement admissible
\emph{isol\'ement}. Cependant, une mol\'ecule polycyclique compl\`ete
est soumise \`a des contraintes \textbf{globales} qui ne se r\'eduisent
pas \`a la composition de contraintes locales.

Plusieurs hypoth\`eses peuvent expliquer les r\'esultats observ\'es~:

\begin{enumerate}
  \item \textbf{Contraintes globales manquantes.} La somme des angles
        autour d'un sommet interne (partag\'e par 3~cycles) doit valoir
        exactement $360^\circ$ pour un plongement plan. Un pentagone
        ($108^\circ$) et un heptagone ($\approx 128{,}57^\circ$) cr\'eent
        localement un exc\`es ou un d\'eficit angulaire. M\^eme si chaque
        voisinage local est admissible, l'accumulation de ces d\'eficits
        sur une mol\'ecule de 16~cycles peut rendre le plongement plan
        impossible.

  \item \textbf{Courbure de Gauss discr\`ete.} Par analogie avec la
        g\'eom\'etrie diff\'erentielle, un pentagone introduit une
        courbure positive et un heptagone une courbure n\'egative.
        Pour que la surface reste plane, les courbures doivent se
        compenser globalement. Cet \'equilibre n'est pas captur\'e par
        des contraintes purement locales.

  \item \textbf{Taille du benz\'eno\"ide.} Le cas \texttt{first.graph}
        (3~hexagones) est suffisamment petit pour qu'une seule
        substitution 5--7 reste plane. \`A~16~hexagones, les d\'eformations
        cumul\'ees deviennent trop importantes.

  \item \textbf{Flexibilit\'e mol\'eculaire.} Il est possible que
        certaines configurations soient g\'eom\'etriquement viables avec
        des cycles \textbf{non r\'eguliers} (angles l\'eg\`erement
        d\'eform\'es), ce que notre reconstruction en polygones r\'eguliers
        ne capture pas.  xTB pourrait trouver un minimum plan si les
        coordonn\'ees initiales \'etaient plus proches de l'\'equilibre.
        Cette hypoth\`ese semble n\'eanmoins peu probable, car xTB est
        capable de corriger des g\'eom\'etries assez \'eloign\'ees du
        minimum.
\end{enumerate}

\subsection{Pistes de travail}

\begin{enumerate}
  \item \textbf{Validation par le chimiste.} Soumettre les 41~structures
        XYZ et les r\'esultats xTB au chimiste du projet pour analyse.
        Le chimiste pourra d\'eterminer si certaines configurations sont
        physiquement impossibles ou si des ajustements au mod\`ele
        sont n\'ecessaires.

  \item \textbf{Contrainte globale de courbure.} Explorer l'ajout d'une
        contrainte CSP portant sur la somme des d\'eficits angulaires~:
        $\sum_v (6 - x_v) \cdot \delta_v = 0$, o\`u $\delta_v$
        repr\'esente le d\'eficit angulaire li\'e au changement de taille
        du cycle~$v$. Une telle contrainte \'eliminerait les
        configurations qui introduisent une courbure globale non nulle.

  \item \textbf{Test sur des benz\'eno\"ides interm\'ediaires.}
        Tester avec des benz\'eno\"ides de taille~4 \`a~10 hexagones pour
        identifier \`a partir de quelle taille les solutions cessent
        d'\^etre planes.

  \item \textbf{Exploration de la g\'eom\'etrie initiale.} Si
        l'hypoth\`ese~4 est jug\'ee plausible, explorer des m\'ethodes
        de placement initial plus sophistiqu\'ees (relaxation
        it\'erative, placement par ressorts) pour fournir \`a xTB des
        coordonn\'ees de d\'epart plus proches de l'\'equilibre.
\end{enumerate}


% ════════════════════════════════════════════════════════════════════
\section*{R\'esum\'e}
% ════════════════════════════════════════════════════════════════════

\begin{notebox}[Bilan]
\textbf{Ce qui fonctionne~:}
\begin{itemize}[nosep]
  \item Le pipeline CSP \'enum\`ere correctement les solutions
        (2~pour \texttt{first}, 41~pour \texttt{second}).
  \item Le pipeline de reconstruction produit des mol\'ecules
        topologiquement correctes avec une g\'eom\'etrie raisonnable.
  \item La validation compl\`ete (reconstruction $\to$ xTB $\to$
        planarit\'e) fonctionne de bout en bout.
\end{itemize}

\textbf{Ce qui reste ouvert~:}
\begin{itemize}[nosep]
  \item Les contraintes locales (table de voisinage) ne suffisent pas
        \`a garantir la planarit\'e globale pour les benz\'eno\"ides de
        grande taille.
  \item Il faut d\'eterminer si des contraintes globales suppl\'ementaires
        peuvent \^etre ajout\'ees au mod\`ele CSP, ou si le probl\`eme
        de planarit\'e doit \^etre trait\'e comme un filtre
        a~posteriori.
\end{itemize}
\end{notebox}

\end{document}
